\chapter{Sestao, 1992: El pueblo y su geografía}

\section{La margen izquierda del Nervión}

Sestao en 1992 es un laberinto de cuestas empinadas, edificios obreros de ladrillo rojo y naves industriales gigantescas encajadas contra la orilla de la ría. La atmósfera está impregnada de un perpetuo olor a óxido, carbón y humedad salobre.

\kobfigure{02_escondite_huerta.jpg}{Las huertas y almacenes abandonados a las afueras de Sestao: el refugio del grupo.}

\section{Lugares clave para la aventura}

\subsection{La Sala de Juegos de Txetxu (Gran Vía)}
El auténtico santuario de los adolescentes del pueblo. Un local ruidoso y saturado de humo donde convergen chavales en chándal de táctel. Las máquinas de \emph{Street Fighter II}, \emph{Tetris} y los futbolines de madera son el epicentro de los rumores locales. Txetxu, el dueño, suele mirar para otro lado si alguien pasa la tarde sin gastar demasiado mientras no arme bulla.

\subsection{El Parque de Markonzaga}
Un parque elevado con bancos de madera carcomida, columpios oxidados y una vista directa hacia los muros exteriores de Babcock \& Wilcox. Es el punto de encuentro natural tras salir de clase del Instituto Ángela Figuera o del Colegio Las Llanas.

\subsection{Babcock \& Wilcox (B\&W) y la Vega de Galindo}
Un inmenso complejo cercado por vallas de alambre de espino en la confluencia del río Galindo. B\&W diseñó calderas pesadas y reactores nucleares en las décadas anteriores. Cuenta con vías de tren propias, depósitos subterráneos de residuos y garitas de seguridad vigiladas por guardias jurados de Prosegur con impermeables y perros adiestrados.

\subsection{La Sociedad de Ciencias Naturales (SCN)}
Ubicada en un edificio sobrio cercano al ayuntamiento (en la antigua sede de la AISS). Fundada en 1920, custodia vitrinas de madera repletas de fósiles del Cretácico, lupas binoculares y detallados mapas topográficos de las cuevas mineras de las Encartaciones y La Arboleda. Aquí trabaja Francisco, rodeado de cajas de amonites y viejos informes geológicos.

\subsection{Los Montes de La Arboleda y el Funicular de La Reineta}
Subiendo en el histórico funicular de madera de La Reineta se accede a los antiguos montes de hierro. Lo que antaño fueron tajos abiertos de extracción son ahora lagos de aguas verdosas rodeados de escombreras de roca y bocaminas abandonadas con carteles oxidados de \emph{«Peligro de derrumbe»}.

\subsection{La Casa del Indiano}
Un palacete señorial decadentista de principios del siglo XX situado entre las calles Txabarri y Miguel de Unamuno. Dos altas palmeras flanquean la verja carcomida. Aunque oficialmente lleva décadas cerrado, los chavales rumorean que de noche se ven luces verdes parpadeando en la buhardilla y que las brújulas giran sin control al pasar cerca de su muro.

\subsection{El Gazteleku}
Antigua sede juvenil municipal, frecuentada por los aficionados a los cómics de \emph{Zitrola}, las revistas de videojuegos y las partidas de rol. Un lugar idóneo para contrastar teorías disparatadas sin que los adultos se rían de ellos.
